\textbf{Bounds Consistency}, briefly \textbf{BC}, is a lower level of consistency than GAC or arc-consistency. Usually, it's used when all the domain values 
are ordered, since it relaxes the domain of any decision variable $X_i$ from $D(X_i)$ to $[min(X_i), \dots, max(X_i)]$. So by relaxing we take in account 
the whole interval, limited by the minimum and maximum values coming from the original domain.
\begin{example}
    e.g. Relaxing a decision variable domain. 
    $$ D(X_i) = \{1, 3, 5\} \rightarrow [1\dots5] $$ 
\end{example}
In BC, we don't talk anymore about \textbf{support} but instead about \textbf{bound support}, since values don't come anymore from a precise domain but 
from a more large range. 
\begin{example}[title = Formal example]
    Any constraint $C$ defined on $k$ decision variables, $C(X_1, \dots, X_k)$, gives us the following set of allowed combinations of values:
    $$ C \subseteq D(X_1) \times \dots \times D(X_k) $$

    Each allowed tuple $(d_1, \dots, d_k) \in C$ where $d_i \in [min(X_i), \dots, max(X_i)]$ is a \textbf{bound support} for $C$. \vspace{7pt}

    Therefore, we say that the general constraint $C(X_1, \dots, X_k)$ is BC if and only if:
    $$ \forall X_i \in \{X_1, \dots, X_k\}, min(X_i) \text{ and } max(X_i) \text{ belong to a bound support} $$
\end{example}
Previously, the GAC requirement was totally different, any value $v \in D(X_i)$ must belong to a whatever \textbf{support}, but now we care only about the
\textbf{minimum} and the \textbf{maximum} values. By BC we make the search fastest than before, because we need to check only two values and not each of them. \vspace{7pt}

However, everything comes with advantages and disadvantages, in particular:
\begin{itemize}
    \renewcommand{\labelitemi}{-}
    \item \textbf{Disadvantage}.
    \begin{enumerate}
        \item BC might not detect all GAC inconsistencies. BC checks only the minimum and maximum values, between them might be different undected inconsistencies.
    \end{enumerate}
    \item \textbf{Advantages}.
    \begin{enumerate}
        \item Might be easier to look for a support in a range than in a domain, in particular when domains are very large.
        \item Achieving BC is enough to achieve also GAC when we are dealing with \textbf{monotonic constraints}.
        \item Achieving BC is often cheaper than achieving GAC.
    \end{enumerate}
\end{itemize}
\begin{example}
    e.g. Differences between GAC and BC regarding $\textit{alldifferent}(\dots)$ constraint.
    $$ D(X_1) = D(X_2) = D(X_3) = \{1, 3\}, C:\textit{alldifferent}([X_1, X_2, X_3]) $$
    It's quite easy to see that the constraint $C$ is unsatisfiable, the key distinction lies in when GAC and BC actually recognize this impossibility. \vspace{7pt}
    
    GAC would be perfectly aware that there are three different decision variables and only two available values in the domain and it will just fail.
    BC won't see that, since it uses a range instead of the original domain, assuming also $2$ as available value for the future assignments. Anyway, also BC
    will discover that this constraint is unfeasible, but this will occur later than by GAC. \vspace{7pt}

    Technically, things go as follows:
    \begin{itemize}
        \renewcommand{\labelitemi}{-}
        \item $\text{BC}(C)$.
        \begin{itemize}
            \item All $min(X_i)$ and $max(X_i)$ belong to a bound support.
            \item No domain reduction with BC propagation, since the tuple $d_i$ belongs to the bound support of $C$.
        \end{itemize} 
        \item $\text{GAC}(C)$.
        \begin{itemize}
            \item None of $min(X_i)$ and $max(X_i)$ belongs to a support. 
            \item $C$ fails with GAC propagation.
        \end{itemize} 
    \end{itemize}
\end{example}
\begin{example}
    e.g. Differences between GAC and BC regarding \textbf{monotonic constraint}. 
    $$ D(X_1) = [2\dots6], D(X_2) = [1\dots5], C: X_1 \leq X_2 $$
    In this example, GAC and BC will give the same reasoning, so we can summerize exactly the same behavior as follows:
    \begin{itemize}
        \renewcommand{\labelitemi}{-}
        \item All values $D(X_1)$ must be less or equal to $max(X_2)$.
        \item All values of $D(X_2)$ must be greater or equal to $min(X_1)$.
        \item It's enough to adjust the maximum value of $X_1$ and minimum value of $X_2$ to achieve consistency.
    \end{itemize} \vspace{7pt}
    After that, the final domains will be both $D(X_1) = D(X_2) = [2\dots5]$.
\end{example}